\chapter{Programming in Mathilda}
\label{ch:programming}

Everything in the previous chapters was, in a sense, one long computation: you
typed an expression, Mathilda rewrote it until it stopped changing, and printed
the result. This chapter is about writing \emph{your own} rewrites---about
turning Mathilda from a calculator you drive into a language you program. It
supports three styles, and the point of the chapter is that they are not three
separate languages bolted together but three faces of one mechanism. A
\emph{pattern-based} program describes what an expression looks like and how to
transform it. A \emph{procedural} program says do this, then that, changing
variables as it goes. A \emph{functional} program builds a result by applying
and composing functions rather than by mutating anything. You can mix all three
in a single line, because underneath they are the same thing: rules, applied by
the evaluator to expressions, under the control of a handful of
attributes.\index{programming language}

That last clause is the load-bearing one, and it is worth stating before any
syntax. Two facts from the system's design govern everything here. First,
\emph{everything is an expression}\index{expression!uniformity}---a list, a rule,
a pattern, a function definition, a loop are all trees with a head and
arguments, so the same tools operate on all of them. Second, \emph{attributes
drive evaluation}\index{attribute}: before Mathilda touches a call \texttt{f[\ldots]}
it reads the attributes of \texttt{f}, and those bits decide whether the
arguments are evaluated or held, threaded over lists, flattened, or sorted. Once
you see the paradigms as different ways of installing and firing rules, and the
attributes as the switches that shape when rules fire, the language stops being a
grab-bag of commands and becomes a single coherent machine. We build up to that
machine deliberately, closing the chapter with the attributes that were quietly
running the whole time.

We take pattern matching first, because it is the most characteristic style and
the one the other two lean on. Then procedural programming---scoping, branching,
and loops---for the moments when a plain imperative recipe is the clearest thing
to write. Then functional programming, the shortest and often the fastest code.
And finally evaluation control: how to stop the evaluator, restart it, and hold a
head inert, which is where the attribute story comes home.

\section{Pattern matching and rules}
\label{sec:prog-patterns}

A \emph{pattern}\index{pattern matching} is a template that describes a class of
expressions. Matching a pattern against a subject is the act of unifying the two
structurally---checking that the subject has the described shape and, along the
way, binding names to the pieces that filled the blanks. This one operation is
the engine beneath \B{Cases}, \B{Replace}, every rule you write with
\texttt{->}, and every function you define with \texttt{:=}. Learn it once and
you have learned the spine of the language.

\subsection{The blank and its kin}
\label{ssec:prog-blank}

The atom of every pattern is the \emph{blank}, written \texttt{\_}
(\B{Blank}\index{pattern matching!Blank (underscore)}): it matches any single
expression. A blank can be \emph{typed} by writing a head after it---\texttt{\_Integer}
matches only an expression whose head is \texttt{Integer}---and it can be
\emph{named}, as in \texttt{x\_}, which matches anything and binds the name
\texttt{x} to whatever it matched. Two relatives match \emph{runs} of arguments
rather than a single one: \B{BlankSequence} (\texttt{\_\_}) matches one or more
expressions, and \B{BlankNullSequence} (\texttt{\_\_\_}) matches zero or more.
\B{MatchQ}\texttt{[expr, patt]} asks the yes/no question directly, and is the
quickest way to build an intuition for what fits what.

\begin{codepairs}
\pair{07-programming/patterns/1}{The bare typed blank: \(42\) has head
\texttt{Integer}.}
\pair{07-programming/patterns/2}{\(3.14\) does not---its head is \texttt{Real}, so
\texttt{\_Integer} fails.}
\pair{07-programming/patterns/3}{\ldots but \texttt{\_Real} matches it. A typed
blank is a structural test on the head.}
\pair{07-programming/patterns/4}{Head-typed blanks reach the atomic numeric heads:
\(1/2\) is \texttt{Rational[1, 2]}.}
\pair{07-programming/patterns/5}{And \(3 + 4\,i\) is \texttt{Complex[3, 4]}.}
\pair{07-programming/patterns/6}{Structure, not just heads: \texttt{f[\_, \_]}
matches a two-argument \texttt{f}.}
\pair{07-programming/patterns/7}{\texttt{\_\_} (\B{BlankSequence}) matches a run of
one or more arguments.}
\pair{07-programming/patterns/8}{An empty \texttt{f[]} has no arguments, so
\texttt{\_\_} (one or more) fails\ldots}
\pair{07-programming/patterns/9}{\ldots while \texttt{\_\_\_}
(\B{BlankNullSequence}, zero or more) matches it.}
\pair{07-programming/patterns/10}{Sequence blanks match \emph{positions}: some
prefix, a \(3\), some suffix.}
\pair{07-programming/patterns/11}{A named typed blank binds: the rule fires only on
the integer-argument \texttt{f} and rewrites it, leaving the others untouched.}
\end{codepairs}

The last line is the whole idea in miniature. The named blank \texttt{n\_Integer}
does two jobs at once---it \emph{restricts} the match to an integer argument and
it \emph{captures} that argument as \texttt{n}, ready to be used on the
right-hand side. The two calls that do not fit, \texttt{f[x]} and \texttt{f[1, 2]},
are simply left alone; a pattern that does not match is not an error, it is a
non-event.\index{pattern matching!named pattern}

\calloutnote{Under the hood}{Matching is structural tree unification with
backtracking, in \texttt{src/match.c}. A named pattern binds its variable in a
\texttt{MatchEnv}; a sequence blank (\texttt{\_\_}, \texttt{\_\_\_}) forces a
\emph{search} over how to partition the remaining arguments, and the matcher
backtracks through the partitions until the rest of the pattern also fits. A
non-linear pattern---the same name used twice, as in \texttt{f[x\_, x\_]}---simply
checks that the two bindings are structurally equal, which is why it can express
``two equal arguments'' with no extra syntax.}

\subsection{Selecting by pattern}
\label{ssec:prog-select}

Once a pattern can say \emph{what to look for}, a family of functions puts it to
work searching a larger expression. \B{Cases} keeps the elements that match;
\B{Count} tallies them; \B{Position} reports where they are; \B{DeleteCases}
removes them; and \B{Select}, the workhorse's cousin, keeps elements for which a
\emph{predicate function} returns \texttt{True}. A rule attached to \B{Cases}
transforms each match as it is collected.

\begin{codepairs}
\pair{07-programming/cases/1}{A deliberately mixed bag---integers, a string, a
real, symbols, a power.}
\pair{07-programming/cases/2}{\B{Cases} with \texttt{\_Integer} keeps the three
integers.}
\pair{07-programming/cases/3}{With the predicate pattern \texttt{\_?NumberQ} it
keeps every number, real included.}
\pair{07-programming/cases/4}{\B{Count} returns how many matched rather than which.}
\pair{07-programming/cases/5}{\B{Position} gives the paths of the matches---here the
two places \texttt{b} occurs.}
\pair{07-programming/cases/6}{\B{Select} filters by a predicate: the primes up to
\(20\).}
\pair{07-programming/cases/7}{\B{DeleteCases} is the complement of \B{Cases}: drop
the \texttt{Missing[]} markers.}
\pair{07-programming/cases/8}{A rule on \B{Cases} transforms each match---here each
pair to its sum.}
\end{codepairs}

The division of labor is worth noticing. \B{Cases} and its relatives take a
\emph{pattern} and match structure; \B{Select} takes a \emph{function} and tests
each element. They overlap---\texttt{Cases[list, \_?EvenQ]} and
\texttt{Select[list, EvenQ]} agree---but the pattern form can do things a
predicate cannot, such as destructure (\texttt{\{a\_, b\_\}}) and rewrite in the
same breath, as pair~8 does.

\subsection{Constraining a match}
\label{ssec:prog-constrain}

A bare blank is often too generous. Three constructs narrow it. A
\emph{pattern test}, \texttt{patt\,?\,f}, keeps the match only if the function
\texttt{f} applied to it returns \texttt{True}. A \emph{condition},
\texttt{patt\,/;\,cond}, keeps it only if the boolean expression \texttt{cond}---which
may mention the pattern's bound names---is \texttt{True}. And
\emph{alternatives}, \texttt{p1\,|\,p2\,|\,\ldots}, match any one of several
patterns. The complementary \texttt{Except[c]} matches anything that does
\emph{not} match \texttt{c}.

\begin{codepairs}
\pair{07-programming/constrain/1}{\texttt{\_?EvenQ}: the blank, then the predicate
\B{EvenQ} applied to what it matched.}
\pair{07-programming/constrain/2}{\texttt{n\_ /; n > 7}: a condition on the bound
name \texttt{n}---an arbitrary boolean, not just a head test.}
\pair{07-programming/constrain/3}{\texttt{2 | 4 | 6}: alternatives match any one of
the listed forms.}
\pair{07-programming/constrain/4}{A condition drives a rewrite: negate exactly the
negative elements. This is \(\lvert\cdot\rvert\) written as a rule.}
\pair{07-programming/constrain/5}{A condition may relate two bound names:
\texttt{a < b} keeps only the increasing pairs.}
\pair{07-programming/constrain/6}{\texttt{Except[\_String]} keeps everything that is
not a string.}
\end{codepairs}

The difference between \texttt{?} and \texttt{/;} is the difference between a
test on the \emph{whole match} and a test on its \emph{named parts}.
\texttt{\_?f} hands the matched expression to \texttt{f} and asks for a verdict;
\texttt{patt /; cond} evaluates an arbitrary predicate that can see every name the
pattern bound, so \texttt{\{a\_, b\_\} /; a < b} can compare two pieces at once, a
thing no single-argument predicate test can do. Reach for the test when the
condition is ``this one thing passes \texttt{f}'', and for the condition when the
constraint couples several bindings.

\calloutnote{Theory}{A condition can make matching \emph{undecidable} in general,
because \texttt{cond} is an arbitrary Mathilda computation---the matcher cannot
know whether it halts. This is the price of the expressiveness: the same
generality that lets \texttt{/;} express any side condition means a pattern is a
tiny program, not a finite automaton. In practice conditions are cheap boolean
checks, but it is worth remembering that ``does this match?'' is, with
\texttt{/;} in play, as hard as ``does this program return \texttt{True}?''}

\subsection{Rules and replacement}
\label{ssec:prog-rules}

A \emph{rule} pairs a pattern with a replacement. \B{Rule} (\texttt{lhs -> rhs})
is \emph{immediate}: its right-hand side is evaluated once, when the rule is
built. \B{RuleDelayed} (\texttt{lhs :> rhs}) is \emph{delayed}: the right-hand
side is held and re-evaluated afresh on every match, after the pattern's names
are bound. Rules are applied by three operators that differ only in \emph{where}
and \emph{how often} they look: \B{Replace} tries the rule at one level (the
whole expression by default), \B{ReplaceAll} (\texttt{/.}) sweeps the entire tree
top-down, and \B{ReplaceRepeated} (\texttt{//.}) applies \texttt{/.} over and over
until the expression stops changing.

\begin{codepairs}
\pair{07-programming/rules/1}{\B{Replace} rewrites the whole expression when it
matches the rule's left-hand side.}
\pair{07-programming/rules/2}{\B{ReplaceAll} (\texttt{/.}) finds the \(x^2\)
\emph{inside} the sum and rewrites it in place.}
\pair{07-programming/rules/3}{It rewrites every occurrence in one downward sweep:
both \texttt{x} and the \texttt{x} inside \(x^2\).}
\pair{07-programming/rules/4}{A list of rules is applied together---a simultaneous
substitution.}
\pair{07-programming/rules/5}{Given a list of \emph{rule lists}, \B{Replace} maps
over them, giving one answer per alternative.}
\pair{07-programming/rules/6}{\texttt{/.} applies a rule once per part and does not
re-descend into what it just built---so one \texttt{f} is peeled off.}
\pair{07-programming/rules/7}{\texttt{//.} (\B{ReplaceRepeated}) repeats to a fixed
point, peeling every \texttt{f} until nothing changes.}
\pair{07-programming/rules/8}{\B{Rule} (\texttt{->}) evaluates its RHS \emph{once};
the counter is bumped at rule-build time, so all three parts get the same \(1\).}
\pair{07-programming/rules/9}{\B{RuleDelayed} (\texttt{:>}) re-evaluates its RHS per
match, so the counter climbs \(1, 2, 3\).}
\end{codepairs}

Pairs~6 and~7 draw the line between \texttt{/.} and \texttt{//.}: the single
sweep peels one layer, the repeated sweep runs to a fixed point. Pairs~8 and~9
draw the sharper and more consequential line between \texttt{->} and
\texttt{:>}.\index{rule!immediate vs delayed} With a passive right-hand side the
two are indistinguishable, but the moment the RHS \emph{does} something---calls a
random generator, reads a counter, evaluates a costly function---the choice
matters. Use \texttt{:>} whenever the right-hand side should be recomputed for
each match, which is nearly always the case for rules with bound pattern
variables.

\calloutnote{Pitfall}{\texttt{/.}\ descends into \emph{heads} as well as
arguments, because a head is just another part of the tree. A blank rule like
\texttt{\_ -> new} applied to \texttt{\{a, b\}} will therefore also match the head
\texttt{List} and replace it---rarely what you want. When you mean ``the elements,
not the structure,'' constrain the pattern (\texttt{\_Symbol}, say) or apply the
rule at a definite level with \texttt{Replace[expr, rule, \{1\}]}, as pairs~8 and~9
do to keep the counter demonstration clean.}

\usagebox{ReplaceAll}

\subsection{Functions as rules: DownValues}
\label{ssec:prog-downvalues}

Here the whole apparatus pays off. Defining a function with \texttt{:=}
(\B{SetDelayed}) does not create a new kind of object---it \emph{installs a
delayed rule} on a symbol. \texttt{square[x\_] := x\^{}2} attaches to the symbol
\texttt{square} the rule ``a call matching \texttt{square[x\_]} rewrites to
\texttt{x\^{}2}''. These rules are the symbol's
\emph{DownValues}\index{DownValues}, and every time you write \texttt{square[7]}
the evaluator matches the call against them, exactly as \texttt{/.}\ would. A
function \emph{is} a bundle of pattern rules.

\begin{codepairs}
\pair{07-programming/downvalues/1}{Install a rule on \texttt{square} with
\texttt{:=}.}
\pair{07-programming/downvalues/2}{It fires by matching: on a number, a symbolic
sum, and a constant alike.}
\pair{07-programming/downvalues/3}{A second symbol, \texttt{area}, gets its own
rule\ldots}
\pair{07-programming/downvalues/4}{\ldots and a \emph{second} rule with a different
argument count.}
\pair{07-programming/downvalues/5}{Dispatch is by pattern: one argument gives the
disk area, two the rectangle.}
\pair{07-programming/downvalues/6}{The base cases of \texttt{fib} are installed with
\texttt{=} (immediate).}
\pair{07-programming/downvalues/7}{The second base case.}
\pair{07-programming/downvalues/8}{The recursive case, with \texttt{:=}, refers to
\texttt{fib} of smaller arguments.}
\pair{07-programming/downvalues/9}{\texttt{fib[10]} unwinds through the rules to
\(55\).}
\pair{07-programming/downvalues/10}{\texttt{collatz} on even integers---the guard
\texttt{\_?EvenQ} selects the rule\ldots}
\pair{07-programming/downvalues/11}{\ldots and a separate rule handles the odd case.}
\pair{07-programming/downvalues/12}{\B{NestList} then iterates the two-rule function:
the Collatz orbit of \(27\).}
\end{codepairs}

The \texttt{fib} definition shows why the ordering of rules matters and how
recursion falls out for free. The two base cases \texttt{fib[0]} and
\texttt{fib[1]} are \emph{more specific} than the general \texttt{fib[n\_]}, and
Mathilda tries specific rules before general ones, so the recursion bottoms out
correctly. The \texttt{collatz} pair shows dispatch by \emph{guard}: two rules
with the same shape but disjoint conditions act as a single case-split function.

\calloutnote{Under the hood}{DownValues live on the head symbol in the symbol
table (\texttt{src/symtab.c}) as a linked list, kept in order of
\emph{specificity}---genuinely more-specific patterns ahead of more-general ones,
with insertion order breaking ties---so that \texttt{fib[0]} is tried before
\texttt{fib[n\_]} however the two were entered (\S\ref{sec:int-symtab}). Evaluating
\texttt{square[7]} is then just \texttt{apply\_down\_values}: walk the list, and
the first pattern that matches wins---the same matcher, the same first-match
rule, as \texttt{/.}. There is no separate ``function call'' machinery in the
kernel; a user function and a rewrite rule are one mechanism.}

This unity buys a classic trick almost for free. A rule whose right-hand side
\emph{installs another rule} turns a function into its own cache:

\mtranscript{07-programming/memoize}

The definition \texttt{fib[n\_] := fib[n] = fib[n-1] + fib[n-2]} reads oddly at
first. Each time it fires for a new \texttt{n}, the inner \texttt{fib[n] = \ldots}
is an immediate assignment that installs a \emph{new}, specific DownValue
\texttt{fib[n] -> value}. The next request for that same \texttt{n} matches the
new specific rule instead of re-entering the recursion. Naive \texttt{fib} makes
exponentially many calls; this one makes \(O(n)\), which is why \texttt{fib[100]}
returns a twenty-one-digit integer instantly rather than after the heat death of
the machine.\index{memoization}\index{dynamic programming}

\calloutnote{Performance}{Memoization by self-installing rules is the CAS
idiom for dynamic programming. The cost is memory---each computed value becomes a
permanent DownValue on the symbol---so it is a space-for-time trade, not a free
lunch; a long-running kernel that memoizes over a wide domain will accumulate
rules until the symbol is cleared. The win is asymptotic and often enormous, but
the stored table is real, and \B{Clear} is how you reclaim it.}

\section{Procedural programming}
\label{sec:prog-procedural}

Sometimes the clearest description of a computation is a recipe: set this up, do
that repeatedly, choose between these branches. Mathilda has the full
imperative apparatus, and because assignments and loops are themselves
expressions with attributes, they compose with the other two paradigms rather
than sitting apart from them.

\subsection{Assignment and scoping}
\label{ssec:prog-scoping}

Two assignment operators anchor procedural code: \B{Set} (\texttt{=}) evaluates
its right-hand side \emph{now} and binds the result, while \B{SetDelayed}
(\texttt{:=}) holds the right-hand side and re-evaluates it on each use---the same
immediate/delayed distinction as \texttt{->} versus \texttt{:>}, and for the same
reason. A sequence of steps is joined with \B{CompoundExpression}
(\texttt{;}), which evaluates each in turn and returns the last.

The subtler question is \emph{scoping}: how to introduce a variable that is local
to a block of code, so it neither sees nor disturbs the outside world. Mathilda
offers three constructs, and the differences between them are exactly the
differences a working programmer needs to understand. \B{Module} creates genuinely
\emph{fresh} local variables (lexical scoping). \B{With} performs lexical
\emph{substitution}, replacing a name by a value textually before the body runs.
\B{Block} does \emph{dynamic} scoping, temporarily overriding a symbol's global
value for the duration of the body and restoring it afterward.

\begin{codepairs}
\pair{07-programming/scoping/1}{A global \texttt{x} exists in the session.}
\pair{07-programming/scoping/2}{\B{Module} makes a fresh, private \texttt{x}; inside,
it is \(2\).}
\pair{07-programming/scoping/3}{The global \texttt{x} is untouched---still \(100\).}
\pair{07-programming/scoping/4}{\B{With} substitutes \texttt{x -> 5} into the body
before evaluating it.}
\pair{07-programming/scoping/5}{Because \B{With} substitutes \emph{textually}, it
reaches inside \B{Hold}: \texttt{a} becomes \(2\) even though the body is held.}
\pair{07-programming/scoping/6}{A \B{Module} local can be read and assigned from
inside a nested pure function---here a running counter.}
\pair{07-programming/scoping/7}{A global \texttt{n}\ldots}
\pair{07-programming/scoping/8}{\ldots and a function \texttt{g} that reads it.}
\pair{07-programming/scoping/9}{\B{Block} overrides \texttt{n} \emph{dynamically}:
\texttt{g[]} sees the temporary \(3\), so \(3^2 = 9\).}
\pair{07-programming/scoping/10}{\B{Module} does \emph{not}: its private \texttt{n}
is invisible to \texttt{g}, which still reads the global \(100\).}
\pair{07-programming/scoping/11}{Afterward the global \texttt{n} is restored, and
\texttt{g[]} is \(100\) again.}
\end{codepairs}

Pairs~9 and~10 are the crux, and the single most useful thing to internalize
about scoping. \B{Block} localizes a \emph{value}: for the duration of its body,
the symbol \texttt{n} has a new value, and \emph{any} code that reads
\texttt{n}---including a function \texttt{g} defined elsewhere---sees it. \B{Module}
localizes a \emph{name}: it invents a brand-new symbol so its \texttt{n} cannot
possibly collide with anyone else's, which is exactly why \texttt{g}, whose body
refers to the ordinary global \texttt{n}, does not see the Module's version.
Module is what you want almost always, because it is hygienic; Block is the right
tool precisely when you \emph{do} want to reach into other code's use of a symbol---temporarily
raising \texttt{\$RecursionLimit}, say, or overriding a constant for one
computation.\index{scoping!lexical}\index{scoping!dynamic}

\calloutnote{Under the hood}{\B{Module} implements freshness by renaming: its
local \texttt{x} becomes a uniquely numbered symbol \texttt{x\$nnn} (drawn from
\texttt{\$ModuleNumber}) carrying the \texttt{Temporary} attribute, which is why
it cannot clash with a global \texttt{x} or with another Module's. \B{With} needs
no rename because it substitutes and then vanishes---and because the substitution
happens \emph{before} the body evaluates, it is the one scoping construct that
penetrates a \B{Hold} (pair~5). \B{Block} saves the symbol's current value, installs
the new one, and restores the old value on exit even if the body throws. All
three are \texttt{HoldAll}, so the body is not evaluated until the locals are in
place.}

\usagebox{Module}

\subsection{Branching}
\label{ssec:prog-branching}

Three constructs choose among alternatives. \B{If}\texttt{[cond, t, f]} is the
two-way branch. \B{Which}\texttt{[t1, v1, t2, v2, \ldots]} walks a list of
test/value pairs and returns the value of the first test that is \texttt{True}---a
multi-way branch, with a trailing \texttt{True} acting as the catch-all.
\B{Switch}\texttt{[e, form1, v1, \ldots]} instead \emph{pattern-matches} the
value \texttt{e} against a sequence of forms, so it dispatches on shape rather
than on a chain of booleans, with a trailing \texttt{\_} as the default.

\begin{codepairs}
\pair{07-programming/branching/1}{\B{If} chooses the true branch.}
\pair{07-programming/branching/2}{\B{Which} as a multi-way classifier\ldots}
\pair{07-programming/branching/3}{\ldots dispatched over a list with
\texttt{/@}.}
\pair{07-programming/branching/4}{The sign function, three cases with a \texttt{True}
default.}
\pair{07-programming/branching/5}{Its values on a negative, zero, and positive
input.}
\pair{07-programming/branching/6}{\B{Switch} dispatches on \emph{pattern}, not on a
boolean chain.}
\pair{07-programming/branching/7}{So it separates an integer, a real, and a symbol
by their heads.}
\pair{07-programming/branching/8}{\B{Switch} inside a pure function, mapped over a
list.}
\end{codepairs}

Which and Switch are the same idea approached from the two sides of the chapter:
Which tests \emph{predicates} (procedural), Switch matches \emph{patterns}
(the rewriting style). Choose Switch when the distinction you care about is
structural---a head, a shape---and Which when it is a numeric or logical
condition.

\calloutnote{Under the hood}{\B{If} is \texttt{HoldRest}, \B{Which} is
\texttt{HoldAll}, and \B{Switch} is \texttt{HoldRest}: the branches are held and
only the chosen one is ever evaluated. Without that, \texttt{If[x > 0, expensive,
cheap]} would compute both arms before choosing---the attribute is not an
optimization but a correctness requirement. This is the first place the attribute
system does visible work; \S\ref{sec:prog-attributes} makes the pattern explicit.
The three lower to machine branches inside \B{Compile} as well, so a classifier
used in a numeric loop costs nothing (Chapter~\ref{ch:compilation}).}

\subsection{Loops}
\label{ssec:prog-loops}

For genuine iteration Mathilda has \B{Do} (repeat over a range), \B{While}
(repeat while a test holds), and \B{For} (the C-style initialize/test/increment
loop). \B{Break} exits the enclosing loop early and \B{Continue} skips to its next
iteration; both take effect the instant they are evaluated.

\begin{codepairs}
\pair{07-programming/loops/1}{\B{Do} accumulates \(1 + 2 + \cdots + 100 = 5050\).}
\pair{07-programming/loops/2}{The same shape builds \(6! = 720\).}
\pair{07-programming/loops/3}{\B{While} doubles until it passes \(100\).}
\pair{07-programming/loops/4}{\B{For} sums the first ten squares, \(385\).}
\pair{07-programming/loops/5}{\B{Break} leaves the loop as soon as \texttt{i}
exceeds \(3\), so the last value recorded is \(3\).}
\pair{07-programming/loops/6}{\B{Continue} skips the even \texttt{i}, summing only
the odds: \(1+3+5+7+9 = 25\).}
\end{codepairs}

Each loop here returns a value through an accumulator variable, which is the
honest way to use them: a loop's own return value is \texttt{Null}. That is the
tell that loops are about \emph{effects}---changing \texttt{s}, \texttt{p},
\texttt{n}---not about producing a value, and it is the seam where procedural code
gives way to the functional style of the next section, which produces values
without mutating anything.

\calloutnote{Performance}{A loop whose body is entirely machine arithmetic takes
an automatic fast path (\texttt{src/numloop.c}): it is compiled to a small
double-stack program and run with no \texttt{Expr} allocation per iteration, so
\texttt{Do[s = s + i, \{i, 1, 10\^{}7\}]} is a tight machine loop, not ten million
evaluator steps. The moment the body does something non-numeric---a \B{Print}, a
symbolic head, a call to a user function---the fast path silently declines and the
interpreter runs the body exactly as written, so the answer is always the
interpreter's. An exact-integer counter accumulates in overflow-checked
\texttt{int64} and promotes to a bignum rather than wrapping, which is why
\texttt{Do[p = p\,i, \{i, 1, 6\}]} is exactly \(720\) and would be exactly
\(25!\) at the top end.}

\section{Functional programming}
\label{sec:prog-functional}

Functional code builds a result by applying and composing functions instead of
mutating variables. It is usually the shortest way to say what you mean, and---because
it avoids explicit loops---often the fastest, since the same numeric
fast path that accelerates \B{Do} also drives \B{Map}, \B{Nest}, and \B{Fold}.

\subsection{Pure functions}
\label{ssec:prog-purefunc}

A \emph{pure} (anonymous) function is one written without a name.
\B{Function}\texttt{[x, body]} names its parameter explicitly; the terser
\texttt{body\,\&} leaves the parameters implicit, referring to them by
\B{Slot}: \texttt{\#} or \texttt{\#1} is the first argument, \texttt{\#2} the
second, and \B{SlotSequence} \texttt{\#\#} is all of them at once.\index{pure function}\index{anonymous function}

\begin{codepairs}
\pair{07-programming/purefunc/1}{\texttt{\#\^{}2 \&} is the squaring function;
\texttt{\#} is its argument.}
\pair{07-programming/purefunc/2}{\texttt{\#1} and \texttt{\#2} are the first and
second arguments.}
\pair{07-programming/purefunc/3}{The named form \B{Function}\texttt{[u, u\^{}3]} is
the same idea, spelled out.}
\pair{07-programming/purefunc/4}{A named function of two parameters.}
\pair{07-programming/purefunc/5}{Pure functions are values: pass one straight to
\B{Map}.}
\pair{07-programming/purefunc/6}{And to \B{Select} as a predicate.}
\pair{07-programming/purefunc/7}{\texttt{\#\#} (\B{SlotSequence}) is the whole
argument run, spliced---here twice.}
\pair{07-programming/purefunc/8}{\texttt{\#2} before \texttt{\#1} reorders the
arguments.}
\end{codepairs}

The \texttt{\&} notation is the connective tissue of functional Mathilda: it lets
you write a one-off function exactly where it is needed---as the thing to map, the
predicate to select by, the step to iterate---without the ceremony of naming it.
Named \B{Function} earns its keep when the body is long enough that \texttt{\#1},
\texttt{\#2}, \texttt{\#3} become hard to read, or when you need its third-argument
form to give the function evaluation attributes of its own.

\calloutnote{Under the hood}{A pure function binds its parameters by
\emph{lexical substitution} (\texttt{src/purefunc.c}), not through the global
symbol table: when \texttt{(\#\^{}2 \&)[a + b]} fires, \texttt{\#} is replaced by
\texttt{a + b} in the body before it evaluates. Nested pure functions scope
correctly---an inner \texttt{\&} establishes a fresh set of slots that shadow the
outer ones---but named-parameter \B{Function}s are deliberately \emph{not}
closures: an inner \texttt{Function[e, \ldots]} does not capture an outer
function's parameter. When you need that capture, inject the value with \B{With}.}

\usagebox{Function}

\subsection{Mapping and applying}
\label{ssec:prog-map}

Two operators cover most higher-order work. \B{Map} (\texttt{/@}) applies a
function to each element of an expression. \B{Apply} (\texttt{@@}) does something
subtler and often surprising to newcomers: it replaces the \emph{head} of an
expression. \texttt{Plus @@ \{1, 2, 3, 4\}} turns the list's head \texttt{List}
into \texttt{Plus}, giving \texttt{Plus[1, 2, 3, 4]}, which evaluates to the sum.
\B{MapThread} and \B{Thread} generalize mapping to several lists at once.

\begin{codepairs}
\pair{07-programming/mapapply/1}{\B{Map} wraps each element in \texttt{box}.}
\pair{07-programming/mapapply/2}{\texttt{/@} is the same operator in infix form.}
\pair{07-programming/mapapply/3}{\B{Apply} replaces the head \texttt{List} with
\texttt{Plus}\ldots}
\pair{07-programming/mapapply/4}{\ldots which \texttt{@@} writes infix; the list
becomes a sum.}
\pair{07-programming/mapapply/5}{With a symbolic head the mechanism is bare:
\texttt{List} becomes \texttt{f}.}
\pair{07-programming/mapapply/6}{\texttt{@@@} applies at level one---the head of each
sublist.}
\pair{07-programming/mapapply/7}{\B{MapThread} zips several lists through a function
of several arguments.}
\pair{07-programming/mapapply/8}{\B{Thread} spreads a function over the list
argument, copying the scalar \texttt{u} across.}
\pair{07-programming/mapapply/9}{\B{Thread} over \texttt{==} turns one list equation
into a list of equations.}
\end{codepairs}

\B{Map} changes the \emph{elements}, \B{Apply} changes the \emph{head}: holding
that distinction straight is most of what it takes to read functional Mathilda
fluently. The two combine constantly---\texttt{Total} is essentially
\texttt{Plus @@}, and a great many one-liners are a \texttt{Map} to transform
elements followed by an \texttt{Apply} to combine them.

\subsection{Building by iteration}
\label{ssec:prog-iterate}

Where a procedural program loops, a functional one iterates a function. \B{Nest}
applies a function \(n\) times; \B{NestList} records every intermediate result.
\B{Fold} threads a binary function through a list, carrying an accumulator;
\B{FoldList} keeps the running history---which is exactly a cumulative reduction.
\B{Composition} (\texttt{@*}) builds the function that applies several functions
in succession.

\begin{codepairs}
\pair{07-programming/iterate/1}{\B{Nest} applies \texttt{f} three times,
symbolically.}
\pair{07-programming/iterate/2}{\B{NestList} keeps every stage of the same
iteration.}
\pair{07-programming/iterate/3}{With a numeric function it collapses to a number:
\((1+\cdot)^2\) thrice from \(1\).}
\pair{07-programming/iterate/4}{A trivial counter, to read the shape of the
output.}
\pair{07-programming/iterate/5}{Newton's iteration for \(\sqrt 2\)---each step
\((x + 2/x)/2\)---converging in a handful of steps.}
\pair{07-programming/iterate/6}{\B{Fold} sums a list through an accumulator, giving
\(10\).}
\pair{07-programming/iterate/7}{\B{FoldList} keeps the partial sums---a cumulative
total.}
\pair{07-programming/iterate/8}{\B{Fold} by Horner's rule assembles the digits
\(1, 9, 2, 6\) into the base-ten number \(1926\).}
\pair{07-programming/iterate/9}{\B{Composition} applies its functions
innermost-first.}
\pair{07-programming/iterate/10}{\texttt{@*} is the same, in infix.}
\end{codepairs}

The Newton example (pair~5) is worth dwelling on, because it shows functional
iteration doing real numerical work: each application of the pure function
\texttt{(\# + 2/\#)/2 \&} is one Newton step for \(\sqrt 2\), and \B{NestList}
exhibits the quadratic convergence---the digits roughly double each step until the
machine-precision limit is reached and the value stops moving. The same
computation as a \B{While} loop would be several lines and a mutable variable;
here it is one expression, and it takes the numeric fast path automatically.

\calloutnote{Performance}{\B{Nest}, \B{NestList}, \B{Fold}, \B{FoldList}, \B{Map}
and their relatives share the loop fast path (\texttt{src/numloop.c}): when the
iterated function is machine arithmetic throughout, the whole iteration runs as
compiled \texttt{double} bytecode with no per-step allocation. And crucially, the
same functional heads \emph{also} lower inside \B{Compile}, so a \B{Fold} or
\B{Nest} written at the prompt is the same \B{Fold} or \B{Nest} that runs inside a
compiled kernel (Chapter~\ref{ch:compilation}). Exactness is preserved position
by position: a value the iteration merely passes through keeps its exact type
even when computed neighbors have gone inexact.}

\section{Controlling evaluation}
\label{sec:prog-evaluation}

All three paradigms rest on the evaluator, and to program well you sometimes need
to \emph{stop} it: to keep an expression in the form you typed rather than the
form it would reduce to. Mathilda gives you fine control over evaluation, and the
tools reveal the machinery that has been running underneath the whole chapter.

\subsection{Holding and releasing}
\label{ssec:prog-hold}

\B{Hold}\texttt{[expr]} keeps \texttt{expr} unevaluated. \B{HoldForm} does the
same but prints without the wrapper showing, so held expressions display
naturally. \B{Evaluate}, inside a held argument, forces \emph{that} argument to
evaluate anyway. And \B{ReleaseHold} strips the wrappers, letting the held
expression finally run.

\begin{codepairs}
\pair{07-programming/holds/1}{\B{Hold} keeps \(1+1\) from evaluating.}
\pair{07-programming/holds/2}{\B{HoldForm} holds too, but prints without showing the
wrapper.}
\pair{07-programming/holds/3}{\B{Hold} shields all of its arguments.}
\pair{07-programming/holds/4}{\B{Evaluate} punches through the hold for one
argument---\(1+1\) becomes \(2\), the other stays.}
\pair{07-programming/holds/5}{\B{ReleaseHold} removes the wrapper and lets
evaluation resume.}
\pair{07-programming/holds/6}{\B{Hold} holds because it \emph{carries the attribute}
\texttt{HoldAll}.}
\pair{07-programming/holds/7}{So does \B{SetDelayed}---which is why \texttt{:=} can
capture an unevaluated right-hand side.}
\end{codepairs}

Pairs~6 and~7 are the reveal. \B{Hold} is not special-cased in the evaluator; it
holds its arguments because it has the attribute \texttt{HoldAll}, and
\emph{anything} with that attribute holds the same way. \B{SetDelayed} has it too,
which is the whole reason \texttt{f[x\_] := body} can store \texttt{body}
unevaluated instead of computing it once at definition time. Holding is not a
feature of a few privileged heads---it is an attribute, and the next subsection
shows how to make it, and the others, do work.

\subsection{Inactive heads}
\label{ssec:prog-inactive}

Holding an \emph{argument} is one thing; sometimes you want to hold a
\emph{head}---to build \texttt{Plus[2, 3]} or \texttt{Integrate[f, x]} as an inert
object that will not fire until you say so. \B{Inactive}\texttt{[f]} is that inert
form of the head \texttt{f}: \texttt{Inactive[f][args]} evaluates its arguments
but does not run \texttt{f}'s own rules, so the call stays symbolic.
\B{Activate} reverses it, turning every \texttt{Inactive[h]} back into \texttt{h}
and letting evaluation proceed.

\begin{codepairs}
\pair{07-programming/inactive/1}{\B{Inactive}\texttt{[Plus]} is an inert head:
\texttt{Plus} does not fire, so the sum stays symbolic.}
\pair{07-programming/inactive/2}{The \emph{arguments} still evaluate---\(1+1\) and
\(2+2\) reduce---but the outer \texttt{Plus} does not.}
\pair{07-programming/inactive/3}{\B{Activate} releases the head and the addition
finally happens.}
\pair{07-programming/inactive/4}{The chief use: holding an integral inert, so no
integration cascade runs.}
\pair{07-programming/inactive/5}{\B{D} applies the fundamental theorem of calculus
to the inert integral \emph{without} evaluating it.}
\pair{07-programming/inactive/6}{And \B{Activate} runs the held integral when you
want the answer.}
\end{codepairs}

The distinction in pairs~1 and~2 is precise and worth stating: \B{Inactive} freezes
the \emph{head}, not the arguments, so \texttt{Inactive[Plus][1+1, 2+2]} evaluates
its arguments to \texttt{Inactive[Plus][2, 4]} and stops there. This is more than
a curiosity. An inert \texttt{Inactive[Integrate][g, x]} lets you write down an
integral you do not (yet) want to compute---perhaps because the integrand is
non-elementary and the attempt would be expensive---and still manipulate it
symbolically. Pair~5 shows the payoff: \B{D} knows that differentiating an integral
gives back the integrand, and applies that rule to the inert form directly, so an
inert first integral can be verified without ever paying for the
integration.\index{Inactive!fundamental theorem of calculus}

\calloutnote{Under the hood}{Inertness is automatic, not a special rule.
\texttt{Inactive[f][args]} is an expression whose head is the compound
\texttt{Inactive[f]}, and that compound head carries no DownValues, so the call is
already a fixed point of the evaluator---there is nothing to fire. It is the same
structural trick as \texttt{Derivative[n][f][x]}, a call through a compound head.
\B{Activate} is then just a targeted replacement of \texttt{Inactive[h]} by
\texttt{h} followed by re-evaluation.}

\subsection{Attributes drive it all}
\label{sec:prog-attributes}

We can now name the mechanism that has been at work in every section. Before
Mathilda processes a call \texttt{f[a1, \ldots, aN]}, it reads the
\emph{attributes} of \texttt{f} and lets them shape the step: which arguments to
evaluate (the \texttt{Hold} family), whether to thread over lists
(\texttt{Listable}), whether to flatten nested same-head calls (\texttt{Flat}),
whether to sort the arguments (\texttt{Orderless}). \B{Attributes} reports them,
\B{SetAttributes} adds them, and \B{ClearAttributes} removes them---and you can set
them on your own symbols, watching the evaluator's behavior change.

\begin{codepairs}
\pair{07-programming/attributes/1}{\B{Plus} is \texttt{Flat} (associative),
\texttt{Listable}, \texttt{Orderless} (commutative), and more---its algebra is its
attributes.}
\pair{07-programming/attributes/2}{Give a fresh symbol \texttt{circ} the attribute
\texttt{Orderless}\ldots}
\pair{07-programming/attributes/3}{\ldots and its arguments are sorted into canonical
order automatically.}
\pair{07-programming/attributes/4}{Make \texttt{gp} \texttt{Flat}\ldots}
\pair{07-programming/attributes/5}{\ldots and nested \texttt{gp} calls flatten into
one.}
\pair{07-programming/attributes/6}{Make \texttt{hl} \texttt{Listable}\ldots}
\pair{07-programming/attributes/7}{\ldots and it threads element-wise over a list
argument, with no loop written.}
\pair{07-programming/attributes/8}{Over several list arguments it threads them in
parallel.}
\pair{07-programming/attributes/9}{Remove the attribute\ldots}
\pair{07-programming/attributes/10}{\ldots and the threading stops---\texttt{hl}
holds the list whole again.}
\end{codepairs}

This is the design at the heart of Mathilda, and it is why the language feels so
uniform. The evaluator itself is small and generic; it does not know that
\B{Plus} is commutative or that \B{Sin} threads over lists. It knows only how to
read attribute bits and act on them. Commutativity is \texttt{Orderless};
associativity is \texttt{Flat}; element-wise threading is \texttt{Listable};
holding an argument is one of the \texttt{Hold} bits. A new builtin opts into any
of these behaviors by setting a flag, and the same seven or eight attributes
compose to describe everything from arithmetic to your own functions. The three
paradigms of this chapter are three ways of installing rules; the attributes are
the dial that sets \emph{when and how} those rules fire.\index{attribute!Flat}\index{attribute!Orderless}\index{attribute!Listable}\index{attribute!HoldAll}

\calloutnote{Theory}{The evaluator's per-call sequence is fixed and worth
memorizing: evaluate the head; read its attributes; evaluate the arguments except
those held; thread if \texttt{Listable}; flatten if \texttt{Flat}; sort if
\texttt{Orderless}; then try the builtin, then the DownValues. Each attribute is a
step in that pipeline, which is why \texttt{Orderless} on \texttt{circ} and
\texttt{Flat} on \texttt{gp} take effect the instant they are set, with no change
to the symbol's rules---they change how the evaluator \emph{prepares} the call
before any rule is consulted. The mechanism lives in \texttt{src/eval.c} and
\texttt{src/attr.c}.}

\section*{Where this connects}

Programming in Mathilda is not a layer on top of the algebra system; it is the
same system, used deliberately. A function you define with \texttt{:=} is a rule
in the symbol table, matched by the same engine that powers \B{Solve}'s
case analysis and the derivative rules of Chapter~\ref{ch:mathematics}; the
pattern language of \S\ref{sec:prog-patterns} is the language those internal rule
sets are written in. The functional iteration of \S\ref{ssec:prog-iterate} and the
loops of \S\ref{ssec:prog-loops} share one numeric fast path, and both lower into
the bytecode compiler---so the \B{Map}, \B{Fold}, and \B{Nest} you write here are
exactly what \B{Compile} accelerates in Chapter~\ref{ch:compilation}, where the
compilable subset and its cliff are the subject. The attributes of
\S\ref{sec:prog-attributes} are the same bits Chapter~\ref{ch:internals} traces
through the evaluator, and the ``everything is an expression'' model that lets
\B{Map} and \B{ReplaceAll} work on rules, patterns, and functions alike is the
data model of Chapter~\ref{ch:datastructures}. As always, the exhaustive options
of any function named here are one \texttt{?Name} away at the prompt.
